% Source: https://artofproblemsolving.com/community/c3h2755335_a_conjecture
% Stable identity: aops-topic-2755335-q01
If $m=1$, then $n^n=1$, so $n=1$; if $n=1$, then $m^{m^m}=1$, so $m=1$. Assume from now on that $m,n\ge2$.

Since $n^n=m^{m^m}$, the numbers $m$ and $n$ have the same prime divisors. For each such prime $p$, comparing exponents of $p$ gives $n\,v_p(n)=m^m\,v_p(m)$, so
\[\dfrac{v_p(n)}{v_p(m)}=\dfrac{m^m}{n}=:x\]
is the same positive rational number for every $p$. Hence $n=m^x$, and then
\[x=\dfrac{m^m}{n}=m^{m-x}.\]
\begin{itemize}
\item If $x\ge m$, then $m^{m-x}\le1$, so $x\le1<2\le m$, a contradiction.
\item If $x<m$, write $m-x=\dfrac uv$ in lowest terms with $u,v\ge1$. Then $x^v=m^u$ is an integer, so the rational number $x$ is an integer. Every prime exponent of $m^u$ is a multiple of $v$; since $\gcd(u,v)=1$, every prime exponent of $m$ is a multiple of $v$, so $m=z^v$ for an integer $z$. As $x$ and $m$ are integers, $m-x=\dfrac uv$ is an integer, so $v=1$ and $m-x\ge1$. But then $x=m^{m-x}\ge m$, a contradiction.
\end{itemize}
So no solution has $m,n\ge2$, and the only solution is $\boxed{m=n=1}$.
